\chapter*{阅读指引}

一批历史邮件可以提供类别答案，却不会自动给出判断新邮件的规则。我们还要决定读取哪些信息、
允许哪些计算关系、怎样根据答案选择参数，以及如何用未参与选择的邮件检验结果。本部分从
这一过程建立任务、经验、模型、学习和评价的含义，为后续各卷提供共同的讨论对象。

《引言》沿垃圾邮件分类的准备、训练、预测和评价展开：每一步都说明使用什么信息、得到什么
结果，并区分训练误差与新邮件上的错误。《模型简介》继续固定任务，比较权重、条件规则和
实例等表示，再从规模、概率、关系、输出及学习准则解释模型的差别。《范式简介》改变经验
条件，讨论标签不足、行动只有后果、新旧任务和分散数据出现时，学习怎样组织。

读完三章，我们应能说明一个系统要预测或执行什么、模型保留了哪些关系、训练依据来自哪里，
以及现有评价能够支持多大范围的结论。理论部分随后固定抽样方式、损失和比较基准，研究
有限经验何时足以支持学习保证。可信性部分再检查保证之外的使用条件：预测是否可靠、变化
和不同人群怎样影响结果、信息与操作权限怎样受限，以及失效以后如何发现和处置。
